\documentclass[12pt,letterpaper]{article}
\usepackage[utf8]{inputenc}
\usepackage[T1]{fontenc}
\usepackage{lmodern,amsmath,amssymb,array,enumitem,xcolor,geometry,tcolorbox}
\geometry{left=1.7cm,right=1.7cm,top=1.5cm,bottom=1.6cm}
\setlength{\parindent}{0pt}
\definecolor{azul}{HTML}{17365D}
\definecolor{claro}{HTML}{EAF1F8}
\newcommand{\familia}[2]{\section*{Familia #1. #2}}
\newcommand{\campo}[1]{\textcolor{azul}{\textbf{#1}}}
\begin{document}
\begin{center}
{\LARGE\bfseries\color{azul} SUCESIONES Y LÍMITES}\par
{\Large Banco paramétrico -- Nivel básico}
\end{center}
\begin{tcolorbox}[colback=claro,colframe=azul]
Estas familias son originales. Los parámetros se eligen desde las tuplas autorizadas;
no deben combinarse libremente cuando se indique dependencia. El nivel básico privilegia
una técnica por pregunta, índices pequeños y resultados enteros o fracciones simples.
\end{tcolorbox}

\familia{1}{Términos de una sucesión explícita}
\campo{Plantilla:}
\[
 a_n=pn+q.
\]
Pedir tres términos consecutivos y $a_k$.\par
\campo{Condiciones:} $p\neq0$, $1\leq k\leq12$ y $|a_k|\leq40$.\par
\campo{Tuplas recomendadas $(p,q,k)$:}
\[
 (2,1,8),\ (3,-2,7),\ (-2,9,6),\ (4,-3,5).
\]
\campo{Control:} $a_k=pk+q$.

\familia{2}{Reconocimiento de patrón}
\campo{Plantilla:}
\[
 a_n=\frac{n+r}{n+s}.
\]
Mostrar los primeros cuatro términos y pedir el término general y $a_k$.\par
\campo{Condiciones:} $r,s\in\mathbb Z$, $s\geq1$, $r\neq s$, $5\leq k\leq12$.\par
\campo{Tuplas recomendadas $(r,s,k)$:}
\[
 (1,3,8),\ (2,5,7),\ (3,4,10),\ (-1,2,9).
\]
\campo{Control:} mantener denominadores positivos y menores que $16$.

\familia{3}{De forma explícita a recursiva}
\campo{Plantilla:}
\[
 b_n=a_1+(n-1)d.
\]
Pedir $b_1$, la recurrencia y un término posterior.\par
\campo{Condiciones:} $a_1,d\in\mathbb Z$, $d\neq0$, $|d|\leq5$; usar
$b_{n+1}=b_n+d$.\par
\campo{Tuplas recomendadas $(a_1,d,k)$:}
\[
 (3,2,9),\ (7,-1,10),\ (-2,3,8),\ (10,-2,7).
\]
\campo{Control:} $b_k=a_1+(k-1)d$.

\familia{4}{Monotonía inmediata}
\campo{Plantilla:}
\[
 c_n=pn+q.
\]
Pedir decidir si la sucesión es creciente o decreciente mediante
$c_{n+1}-c_n$.\par
\campo{Condiciones:} $p\in\{-4,-3,-2,2,3,4\}$ y $|q|\leq6$.\par
\campo{Regla de control:}
\[
 c_{n+1}-c_n=p.
\]
Si $p>0$ es estrictamente creciente; si $p<0$ es estrictamente decreciente.

\familia{5}{Acotación elemental}
\campo{Plantilla:}
\[
 d_n=L-\frac{m}{n+r}.
\]
Pedir demostrar una cota superior y encontrar el menor término.\par
\campo{Condiciones:} $L,m,r\in\mathbb N$, $1\leq m\leq6$, $0\leq r\leq3$ y
$m$ múltiplo de $r+1$ cuando se desee un primer término entero.\par
\campo{Tuplas recomendadas $(L,m,r)$:}
\[
 (4,2,1),\ (5,3,2),\ (6,4,1),\ (3,2,0).
\]
\campo{Control:} $d_1=L-m/(r+1)$ y $d_n<L$.

\familia{6}{Límite algebraico directo}
\campo{Plantilla:}
\[
 e_n=\frac{an+b}{cn+d}.
\]
Pedir calcular el límite dividiendo por $n$.\par
\campo{Condiciones:} $a,c\in\mathbb N$, $c\neq0$, $|b|,|d|\leq5$,
$cn+d>0$ para $n\geq1$ y $a/c$ entero o fracción irreducible sencilla.\par
\campo{Tuplas recomendadas $(a,b,c,d)$:}
\[
 (2,1,1,3),\ (3,-1,2,2),\ (4,3,2,1),\ (5,-2,3,4).
\]
\campo{Control:} $\displaystyle\lim e_n=a/c$.

\familia{7}{Recurrencia convergente guiada}
\campo{Plantilla:}
\[
 x_1=u,\qquad x_{n+1}=\frac{x_n+L}{2}.
\]
Pedir calcular tres términos, conjeturar monotonía y hallar el límite.\par
\campo{Condiciones:} $u,L\in\mathbb Z$, $u\neq L$, $|u-L|\in\{2,4,8\}$.
Si $u<L$, la sucesión es creciente; si $u>L$, es decreciente.\par
\campo{Tuplas recomendadas $(u,L)$:}
\[
 (0,4),\ (1,5),\ (8,4),\ (10,2).
\]
\campo{Control:} límite $L$ y fórmula
$x_n=L+(u-L)2^{-(n-1)}$.

\end{document}
